% \begin{savequote}
% \includegraphics[width=3.5cm]{./images/pallinger-christoph-aass/Anton_32986_v2-AASS-0112mm.jpg}
% 
% Bild: Christoph Pallinger
% \end{savequote}
\def\ChapterCode{ASS}
\chapter
[\CC{[ASS]} Assistenztätigkeiten]
{Assistenztätigkeiten}
% \AasRsN{RS~XIII}
\label{chp:SAN}

\ChapterIntro
{%
~~
}
{%
\begin{description}
  \item[Maintainer:] Diverse
  \item[Autoren:] Diverse
  \item[Reviewer:] --
  \item[Version:] {\DocVersion} ({\DocVersionInfo})
  \item[SHA1:] {\DocGitCommit}
\end{description}

{\TxTeilkapitelAASS}
}


\Text{Special Credits}{%
\noindent Teile des Textes und der Handlungsanleitungen sind
in Zusammenarbeit mit dem \emph{Department für medizinische
Aus- und Weiterbildung -- Medizinische Universität Wien,
Abt. Methodik und Entwicklung},  welches in Zusammenarbeit 
mit dem Klinischen Institut für Krankenhaushygiene
entstanden. Wir danken Ass.-Prof. Dr. Michael Schmidts, Dr.
Alexander Blacky und deren Team für die Unterstützung!
}{

}
{icons/muw.png}
{}
{MUW}









\clearpage
\section{Arbeiten mit  Medikamenten}

\subsection{Allgemeines zu Medikamenten}

\Text
{Wirkstoff, Dosis, Konzentration}
{Bei Medikamenten wird normalerweise die Menge des
Wirkstoffes (z.\,B. in \Mg) angegeben. Manchmal erfolgt die
Angabe auch in anderen Einheiten, wie z.\,B. den
\I{Internationelen Einheiten} (\I{IE}). Bei Ampullen findet sich außerdem auch eine Volumsangabe.
Aus der Angabe der Wirkstoffmenge und des Volumens kann man
die Konzentration berechnen.


\Layout{57}
{}
{}
{}
{\begin{tabulary}{\linewidth}{LRRCL}
\toprule\addlinespace
\noindent\TabHeading{Form} 	& \noindent\TabHeading{Wirkstoffmenge} 	& \noindent\TabHeading{Volumen} 	& \noindent\TabHeading{Konzentration}							& \noindent\TabHeading{Beispiele}
\\\addlinespace\midrule\addlinespace
Ampulle 			& 2\Mg* 						& 20\Ml* 				& 0,1\nicefrac{\Mg}{\Ml} oder 1\nicefrac{\Mg}{10\Ml*}	& L-Adrenalin
\\
Ampulle 			& 5\Mg* 						& 5\Ml* 				& 1\,\nicefrac{\Mg}{\Ml}								& Dormicum{\texttrademark}
\\
Ampulle 			& 15\Mg* 						& 3\Ml* 				& 5\,\nicefrac{\Mg}{\Ml}								& Dormicum{\texttrademark}
\\
Tabletten 			& 500\Mg* 						&						&														& Parkemed{\texttrademark}, Mexalen{\texttrademark}, \ldots
\\

Trockenstechampulle & 4000 IE						&						&														& Heparin
\\\addlinespace\bottomrule\addlinespace
\end{tabulary}}
{Typische Angaben zu Wirkstoffmenge und Konzentration}
{}





% \begin{itemize}
%   \item Ampulle mit Wirkstoff 2\Mg* in 20\Ml* {\dotfill}
%   Konzentration 0,1\nicefrac{\Mg}{\Ml} oder
%   1\nicefrac{\Mg}{10\Ml*}
%   \item Ampulle mit Wirkstoff 5\Mg* in 5\Ml* {\dotfill}
%   Konzentration 1\,\nicefrac{\Mg}{\Ml}
%   \item Ampulle mit Wirkstoff 15\Mg* in 3\Ml* {\dotfill}
%   Konzentration 5\,\nicefrac{\Mg}{\Ml}
%   \item Tabletten mit Wirkstoff 500\Mg* (pro Tablette)
%   \item Trockenstechampulle mit Wirkstoff 4000 IE
% \end{itemize}

Manche Medikamente können auch zur besseren Dosierung
mit \EE{geeigneten Substanzen} (siehe
Packungsbeilage) verdünnt werden:

\begin{itemize}
  \item Ampulle mit Wirkstoff 15\Mg* in 3\Ml* 
  {\dotfill} Konzentration 5\,\nicefrac{\Mg}{\Ml}
  \item Obige Ampulle 
  + 12\Ml* physiologische Kochsalzlösung 
  (\ce{NaCl} 0,9\PC) 
  {\dotfill} Konzentration 1\,\nicefrac{\Mg}{\Ml}
\end{itemize}

\Cave{Medikamente sind oft auch mit \EE{gleichem Namen} in 
unterschiedlichen Konzentrationen erhältlich 
(z.\,B. Dormicum{\texttrademark}, Ketanest{\texttrademark}).}
\Fazit{Beim Zubereiten bzw. Aufziehen von Medikamenten muss 
man sehr genau auf die Wirkstoffmenge, das Volumen und die 
Konzentration achten.}
}
{}
{}
{}
{}


\subsection{Verabreichungsarten}
% \TextSlide{\TxBeschreibung}{
Es gibt verschiedene Arten Medikamente zu verabreichen. Dies
kann entweder \II{invasiv} (verletzend) oder
\II{nicht-invasiv} sein. 
% }{%
% \begin{itemize}
%     \item Invasiv
%     \item Nicht-invasiv
% \end{itemize}
% }
% {}{}{}

\TextSlide{Invasiv}{%
\index{Injektion!intravenöse}%
\index{Injektion!intramuskuläre}%
\index{Injektion!subkutane}%
\index{Injektion!intraossär}%
\index{intravenös}%
\index{intramuskulär}%
\index{intraossär}%
Bei der invasiven Verabreichung wird das Medikament mittels
einer Kanüle in das jeweilige Gewebe bzw. Gefäß dem
Patienten injiziert. Diese Verabreichungsarten sind
i.\,d.\,R.
den Ärzten, besonders dafür geschulten
Sanitätern\VARIANT{REG-AUT}{ (NKV)}, diplomiertem
Krankenpflegepersonal oder Hebammen vorbehalten.
Die gängigen Injektionsarten sind: 

\begin{itemize}
  \item \EE{Intravenöse Injektion} (\II{i.\,v.}): Das
  Medikament wird mittels Kanüle, Venenverweilkanüle
  o.\,ä. direkt in eine \E{Vene} gespritzt. 
  \item \EE{Intramuskuläre Injektion} (\II{i.\,m.}): Das
  Medikament wird in den \E{Muskel} verabreicht.
  \item \EE{Subkutane Injektion} (\II{s.\,c.}): Die
  Substanz wird in die Fettschicht  der Unterhaut (Subkutis)
  eingespritzt.
  \item Eine Besonderheit ist der \EE{intraossäre Zugang}
  (\II{i.\,o.}), bei dem die Substanz mittels einer
  speziellen Kanüle in die Knochenmarkhöhle verabreicht wird.
\end{itemize}



}{
\begin{itemize}
    \item Intravenös (i.\,v.)
    \item Intramuskulär (i.\,m.)
    \item Subkutan (s.\,c.)
    \item Intraossär (i.\,o.)
\end{itemize}
}{

}{}{}

\TextSlide{Nicht invasiv}{
Bei der nicht invasiven Verabreichung wird das Medikament
über eine Schleimhaut oder direkt über die Haut aufgenommen.
Für die \II{orale} Aufnahme muss das Arzneimittel geschluckt
werden. Bei der \I[sublingual]{sublingualen} Methode kommt
die Substanz unter die Zunge des Patienten. Bei der
\I[perkutan]{perkutanen} Verabreichung wird das Medikament
auf die Haut aufgetragen und aufgenommen.
}{\begin{itemize}
    \item Oral
    \item Sublingual
    \item Perkutan
\end{itemize}}{

}{}



% \clearpage
\subsection{Aufziehen eines Medikamentes aus einer Ampulle}

\Text{\TxBeschreibung}{Glasmpullen lassen sich
unterschiedlich öffnen:
Ampullen mit einem \EE{roten Punkt} haben eine
\E{Sollbruchstelle}, der Daumen muss beim Aufbrechen auf dem
roten Punkt liegen. Ampulle mit einem \enquote{\EE{Halsring}}
können nach allen Richtungen aufgebrochen werden.
Beim Arbeiten muss besonderes Augenmerk auf eine \EE{sterile
Arbeitsweise} gelegt werden! Es kann sonst zu Infektionen
kommen.
}{}{}{}{}

\begin{ProcedureStepsNoHeading}
\Text{Material}{
\begin{compactitem}
  \item Spritze, passende Größe
  \item Kanüle(n) 
  \item Ampulle mit Medikament
  \item Trockene Tupfer
  \item Nadelabwurfbehälter
  \item Ggfs. Verschlussstoppel, Permanentstift
\end{compactitem}
}{}{}{}{}
\end{ProcedureStepsNoHeading}

% \begin{multicols}{2}
% \Text{Vorgehen}{
\begin{ProcedureStepsWide}{Vorgehen}
\begin{enumerate}\CorrectionItemizeBox
  \item Verschlossene \StepHeading{Ampulle vorzeigen} bevor sie geöffnet wird.
  \item \StepHeading{Kontrolle} Ablaufdatum und Inhalt: Verfärbung,
  Ausflockung?
  \item \StepHeading{Öffnen} der Ampulle:
  \begin{itemize}
    \item  Glasampulle:

    Ampullenspitze beklopfen, darauf achten, dass sich
    keine Flüssigkeit im Ampullenkopf mehr befindet

    Öffnen der Medikamenten-Ampulle mit Tupfer, roter Punkt
weist zum Körper

    \item Plastikampulle mit
    \enquote*{Dreh-und-Drink}-Verschluss:
    Verschlussknopf abdrehen.

  \end{itemize}
  \item Ampulle abstellen.
  \item \StepHeading{Spitzenverpackung öffnen}: aufschälen, Spritze in
  Verpackung lassen.
  \item \StepHeading{Kanülenverpackung öffnen}: aufschälen 
  \item \StepHeading{Spritze auf Kanüle}
  aufstecken.
  \item \StepHeading{Kanülenabdeckung entfernen} (nicht abdrehen,
  sondern abziehen!)
  \item \StepHeading{Kanüle in Ampulle} einführen

  Beim Einführen der Kanüle in die Ampulle den
    Ampullenrand nicht berühren.
  \item \StepHeading{Aufziehen} des
  Medikamentes.
  \item Verabreichung:
  \ListBreak{enumerate}{Variante \enquote{Verabreichung über
  Venenzugang}:}
    \item Kanüle \EE{händisch abziehen} und in den
    Nadelabwurfbehälter verwerfen -- \Ca{Keinesfalls Kanüle am
    Nadelabwurfbehälter abstreifen (Kontamination)!}
    \item \StepHeading{Entlüften}
    \item Wird das Medikament nicht sofort verwendet, Spritze mit
    Stoppel verschließen und Ampulle ankleben
    \item Spritze übergeben, Ampulle zeigen
\ListBreak{enumerate}{Variante \enquote{zuspritzen in eine
Infusionsflasche}:}
    \item Gummimembran desinfizieren.
    \item Kanüle nicht bis zum Schaft einstechen.
    \item Medikament zuspritzen.
    \item Beschriften der Flasche.
\ListBreak{enumerate}{Variante \enquote{direktes Spritzen}}
    \item \StepHeading{Kanülengröße} erfragen
    \item \StepHeading{Aufziehkanüle} \Ca{händisch} entfernen
    \item \StepHeading{Entlüften}
    \item \StepHeading{Spritzkanüle} aufstecken
    \item Spritze mit Schutzkappe über Kanüle reichen, Ampulle zeigen
\end{enumerate}
\end{ProcedureStepsWide}
% }{}{}{}{}


\begin{PictureSeriesWide}{}{Bilderserie: \EEE{Medikament aus
Ampulle aufziehen} \SourcePicture{Motal}}
\PictureSeriesRowWide{4}
{./images/motal-aass/medi-kontrolle_4807-00800.jpg}
{Kontrolle}
{./images/motal-aass/medi-aufziehen_4815-00800.jpg}
{Aufziehen mit Kanüle}
{./images/motal-aass/medi-kontrollieren_4824-00800.jpg}
{Übergabe}
{./images/motal-aass/medi-verabreichen_4826_CLOSEUP-00800.jpg}
{Verabreichung über Venflon durch befugtes Personal}
\end{PictureSeriesWide}



% \clearpage
\subsection{Vorbereitung einer Infusion}
\label{topic:infusion}

\begin{ProcedureStepsNoHeading}
\Text{Material}{%
\Married{
\begin{compactitem}
    \item Infusion (Flasche, Beutel)
    \item Infusionsbesteck
    \item Aufhängevorrichtung
    \item Alkoholhaltiges Desinfektionsmittel
    \item einige saubere Tupfer
    \item Permanentmarker
\end{compactitem}
}{
% \centering\includegraphics[width=0.5\linewidth]{./images/infusion_tropfkammer.jpg}
}
}
{}
{}
{}
{}
\end{ProcedureStepsNoHeading}

\begin{ProcedureStepsWide}[Beim Arbeiten muss besonderes Augenmerk auf eine sterile
Arbeitsweise gelegt werden!]{Vorgehen}

\begin{enumerate}
  \item Ggfs. \StepHeading{Aufhängevorrichtung} auf Flasche
  aufsetzen 
  \item \StepHeading{Abdeckung} entfernen (Membranschutz)
  \item Ggfs. \StepHeading{Wischdesinfektion} der Gummimembran
  
  Alkoholgetränkter Tupfer kann während weiterer Vorbereitung
    auf Membran belassen werden.
  \item \StepHeading{Infusionsbesteck montieren}
  \begin{enumerate}
    \item \StepHeading{Durchflußregler schließen} (Radklemme)
    \item Einstich des Tropfkammerdorns in Infusionsflasche mit
    geschlossenem Lüftungsventil und geschlossener Radklemme
    \item Dorn der Tropfkammer und Gummimembran der
    Infusionsflasche dürfen nicht berührt werden!
  \end{enumerate}
  \item \StepHeading{Infusionsschlauch füllen}
  \begin{enumerate}
    \item Flasche wenden und hochhalten
    \item Tropfkammerluftventil öffnen
    \item Tropfkammer durch Zu\-sam\-m\-en\-drück\-en bis zur Hälfte füllen
    \item Konusabdeckung des Infusionsschlauchs am anderen
    Ende muss bei der Belüftung nicht entfernt werden
    \item Infusionsflasche und Ende des Infusionsschlauchs
    mit nicht dominanter Hand halten
    \item Radklemme öffnen
    \item Infusionsschlauch entlüften
    \item Radklemme schließen
  \end{enumerate}
\end{enumerate}
\end{ProcedureStepsWide}


\clearpage % TEMP temp clearpage
% \bigskip
\section{Venenverweilkanüle: Assistenz}
\label{topic:venflon}

\TextSlide{\TxBeschreibung}{
% (\enquote{Venflon}): Vorbereitung und Assistenz 
Periphere Venenverweilkanülen werden routinemäßig zur
intravenösen Gabe von Medikamenten bzw. Infusionen
verwendet. Sie bestehen aus einem Kunststoffschlauch, der in
der Vene zu liegen kommt, und aus einem außenstehenden Teil
mit einer Konnektionsstelle für Infusionen, sowie oft einem
Konus mit Rückschlagventil, für die schnelle Gabe von
Medikamenten. Im Kunststoffschlauch liegt eine
Führungskanüle (\I{Mandrin}), welche das Durchstechen der
Haut und der Vene zwecks der Katheterisierung ermöglicht.
Die Kanüle wird nach Einführen ins Blutgefäß entfernt.

Gebräuchliche Marken sind z.\,B. \I{Venflon{\texttrademark}}
oder \I{Braunüle{\texttrademark}}. Venenverweilkanülen gibt
es in verschiedenen Durchmessern, sie sind farbcodiert.
Normalerweise werden \EE{grüne} Venenverweilkanülen bei
Erwachsenen verwendet, bzw. rosafarbene bei schlechten
Venenverhältnissen.
}{
\begin{itemize}
  \item Kunststoffschlauch, der in
der Vene zu liegen kommt
\item 
Konnektionsstelle für Infusionen
\item Konus mit Rückschlagventil
\item Führungskanüle (\I{Mandrin})
\item Verschiedene Größen, farbcodiert.
Standard: \EE{grün}, rosa bei schlechten
Venen
\item Gebräuchliche Marken:
Z.\,B. \I{Venflon{\texttrademark}},
\I{Braunüle{\texttrademark}}
\end{itemize}

}{}{}{}

\begin{figure}[H]
\includegraphics[width=\linewidth]{./images/demaw/venflon_zerlegt-textwidth.jpg}

\Captionof{figure}{Eine periphere Venenverweilkanüle
(Venflon\texttrademark), in ihre Einzelbestandteile zerlegt.
\SourcePicture{DEMAW}}
\end{figure}



\begin{table}[H]\TableBaseModification
\Captionof{table}{Verschiedene Größen von peripheren
Venenverweilkanülen. \Source{DEMAW}}

\begin{tabular}{m{1.5619999cm}cccll}
\toprule\addlinespace
\TabHeading{Farbcode}
&

&
\TabHeading{Gauge}
&

&
\TabHeading{Länge}
&
\TabHeading{Durchfluss}
\\
 		& 				&  			&{\O} \MM 		& \MM  		& [\nicefrac{ml}{min}] 
\\\addlinespace\midrule\addlinespace
blau 	& Sehr dünn		& \Z{22} 	&\Z{0,5 x 0,8} 	& \Z{25} 	& \Z{25} 
\\
rosa 	& dünn			& \Z{20} 	&\Z{0,7 x 1,0} 	& \Z{32} 	& \Z{55} 
\\
grün 	& normal		& \Z{18} 	&\Z{0,9 x 1,2} 	& \Z{40} 	& \Z{90} 
\\
weiß 	& dick			& \Z{17} 	&\Z{1,1 x 1,4} 	& \Z{42} 	& \Z{135} 
\\
grau 	& Sehr dick		& \Z{16} 	&\Z{1,3 x 1,7} 	& \Z{42} 	& \Z{170} 
\\
orange 	& Noch dicker	& \Z{14} 	&\Z{1,6 x 2,1} 	& \Z{42}
& \Z{265} \\\addlinespace\bottomrule
\end{tabular}
\end{table}


\begin{ProcedureStepsNoHeading}
% \clearpage
% \begin{multicols}{2}
\Text{Material}{
\begin{compactitem}
  \item Nierentasse
  \item Staubinde
  \item Venflon:
  \begin{compactitem}
    \item Die Größe (= Farbe) ist vom Durchführenden vorher zu erfragen!
    \item Die Flügel werden hinunter geklappt
    \item Die farbige Lasche des Zuspritzventils in rechten Winkel zur
    Venflonachse stellen
  \end{compactitem}
  \item Einmalverschlusskappe (\enquote{roter
  Stoppel})
  \item Fixierungspflaster und evtl. zusätzliches
  Fixationsmaterial
  \item Trockene Tupfer
  \item Desinfektions-Tupfer
  \item Spülung, wahlweise:
   \begin{itemize}
    \item Spritze mit \Range{5}{10}\Ml* Spülflüssigkeit
  (\ce{NaCl} 0,9\PC*), oder
  \item Infusion
  \end{itemize}
  \item 3-Wege-Hahn (optional)
  \item (Faserschreiber)
  \item Nadelabwurfbehälter
\end{compactitem}
% \end{multicols}
}{}{}{}{}
\end{ProcedureStepsNoHeading}

\begin{ProcedureStepsNoHeading}
% \begin{multicols}{2}
\Text{Assistenz}{
\begin{enumerate}
  \item Ggfs. \StepHeading{3-Wege-Hahn spülen}
  \begin{itemize}
    \item Spritze mit Spülflüssigkeit ansetzen und durchspülen 
    \item Alternativ, bei Verwendung mit Infusion: 3-Wege-Hahn an Infusionsbesteck anschließen und mit Infusionslösung durchspülen
  \end{itemize}
  \item \StepHeading{Nadelabwurfbehälter} in Reichweite des  Durchführenden
  bereitstellen, sodass er später den Mandrin abwerfen kann.
  \item \StepHeading{Stauschlauch} zureichen
  \item \Z{Der Durchführende legt nun den Stauschlauch an, staut, und sucht eine passende Vene für die Punktion.}
  \item \StepHeading{Alko-Tupfer zureichen}: Mit geöffneter Verpackung sodass dieser
  entnommen werden kann
  \item \StepHeading{Venenverweilkanüle zureichen}:   
  An der Schutzkappe halten, sodass der
  Durchführende die Verweilkanüle abziehen kann
  \item \Z{Der Durchführende punktiert nun die Vene.}
  \item \StepHeading{Trockene Tupfer} zureichen
  \item Je nach Situation 3-Wege-Hahn, Stoppel und Spülung oder Infusionsschlauch
  zureichen
  \item \StepHeading{Fixierpflaster} zureichen
\end{enumerate}
% \end{multicols}
}{}{}{}{}
\end{ProcedureStepsNoHeading}

\begin{PictureSeries}{}{Bilderserie: \EEE{Periphere Venenverweilkanüle}}
\PictureSeriesRow{2}
{./images/demaw/Skriptumaegf-img35}
{Anlage einer periph. Venenverweilkanüle \SourcePicture{DEMAW/Blacky}}
{./images/demaw/Skriptumaegf-img36}
{Eine fixierte Venenverweilkanüle mit Infusion
\SourcePicture{DEMAW/Blacky}} {./images/demaw/venflon_zerlegt}
{\SourcePicture{DEMAW/Blacky}}
{empty}
{}
\end{PictureSeries}

\A{Butterfly gehört erklärt}
% TODO Butterfly


\clearpage
% \bigskip
\section{Endotracheale Intubation: Vorbereitung und
Assistenz} 
\label{topic:intubation-assistenz}%

\noindent Bei der endotrachealen Intubation wird ein
Beatmungsschlauch (\I{Tubus}) über den Mund (seltener über die Nase) durch den
Rachen und den Kehlkopf in die Luftröhre eingeführt. Er
ermöglicht die Beatmung des Patienten bei gleichzeitigem
Freihalten des Atemweges und bietet einen Aspirationsschutz.

\begin{ProcedureStepsWide}[Kontrolle vor Beginn
der Intubation nicht vergessen!]{Material}
\begin{enumerate}
    \item Beatmungsbeutel inkl. passender Beatmungsmaske,
    Bakterienfilter und Sauerstoffreservoir
    \item \ce{O2}-Berieselungseinheit und \ce{O2}-Line
    \item Evtl. Beatmungsgerät
    \item Absaugeinheit: \par\Ca{Keine Intubation ohne
    Absaugung!}
    \item Laryngoskop, bestehend aus:
    \begin{compactenum}
        \item Spatel: gebogen oder gerade; Größen 0\,--\,4
        \item Handgriff mit Batterien
    \end{compactenum}
    \Z{Es gibt Warm- und Kaltlichtlaryngoskope.
    Bei Warmlichtlarynkoskopen befindet sich die Lichtquelle
    am Spatel, bei Kaltlichtlaryngoskpen befindet sich die
    Lichtquelle im Handgriff, das Licht wird über Fiberglas
    in den Spatel geleitet.} \Ca{{\SymbolGefahren}} Daher muss
    darauf geachtet werden, dass die Spatel mit dem Griff
    zusammenpassen! Bei Kaltlichtsystemen gibt es ausserdem
    unterschiedliche Steckverbindungen!
    
    \noindent\begin{minipage}{\linewidth}
\includegraphics[width=\linewidth]{./images/pallinger-christoph-aass/Testlunge_32902-AASS-0112mm.jpg}

\Captionof{figure}{Laryngoskopgriff und verschiedene Spatel
\SourcePicture{Ch. Pallinger}}
\end{minipage}
    \item (Endotracheal-)Tubus; verschiedene Größen
    \item Führungsdraht (\T{Mandrin})
    \item Silikonspray
    \item Evtl. Magillzange
    \item Mit Luft gefüllte Spritze zum Cuffen (je
    nach Tubus, i.\,d.\,R. 10\Ml*)
    \item Stethoskop
    \item Beißschutz: Beißkeil oder Guedel-Tubus
    \item Fixationsmaterial (Breite Heftpflaster, Mullbinde)
\end{enumerate}


\end{ProcedureStepsWide}








\Layout{37}{Abbildung}
{}
{}
{./images/pallinger-christoph-aass/Intubation_32939.jpg}
{Zubehör für die endotracheale Intubation: Absaugeinheit
mit Absaugkatheter, Magill-Zange, Silikonspray, Tubus,
Führungsdraht, Blockerspritze (10\Ml*), verschiedene Spatel,
Laryngoskop, Guedel-Tuben, Beißkeil, Beatmungsbeutel mit
Bakterienfilter, Maske und \ce{O2}-Line, Fixationsmaterial,
Stethoskop.}
{Ch.
Pallinger}

\clearpage



% \MarginBoxFilledRounded{Vorbereitung}{}

\begin{ProcedureStepsWide}{Vorbereitung}
\begin{enumerate}
  \item \StepHeading{Spatelgröße} erfragen
  \item \StepHeading{Tubusgröße} erfragen
  \item \StepHeading{Funktionskontrolle} Laryngoskop (Spatel und
  Griff) zusammenbauen, Funktionskontrolle (Licht)
  
  Ebenso müssen alle möglicherweise benötigten Spatel
  (zumindest $\pm$ 1 Größe und die entsprechenden
  gebogenen bzw. geraden Spatel) mit dem
  vorhandenen Laryngoskopgriff geprüft werden (Lichtquelle
  bei Warmlichtgeräten, passende Verbindung bei Kaltlichtgeräten)
  \item Tubusverpackung aufschälen, Tubus aber in der Verpackung belassen
  \item \StepHeading{Cuff-Dichtigkeit} prüfen: test\-weise aufblasen
  \item \StepHeading{Führungsdraht} (Mandrin) mit Silikonspray
  einsprühen\ACh{HART}{0.8.0}{Mandrin auch mit Silikon einsprühen} und
  einführen. Die Spitze muss bis zum Tubusende vorgeschoben vorgeschoben
  werden, darf aber nicht herausragen.
  \item \StepHeading{Cuff} mit Silikonspray einsprühen
  \item \StepHeading{Stethoskop}: Funktion prüfen
  \item \StepHeading{Absaugungebreitschaft} herstellen:
  \begin{enumerate}
    \item Passenden Absaugkatheter wählen
    \item Verpackung teilweise aufschälen um Anschluss
    freizulegen, Katheter in der Verpackung belassen
    \item Katheter anschließen
    \item Gerät einschalten
    \item Funktionskontrolle: Saugstärke?
    Batteriewarnung?
    \item Absaugeinheit muss in Griff\-weite des Assistenten
    positioniert werden!
  \end{enumerate}
  \item \StepHeading{Beatmungsbeutel} zusammenbauen und an
  \ce{O2}-Berieselungseinheit anschließen
  \item Restliches Material bereitlegen (z.\,B. in Nierentasse)
  \ListBreakHard[Für nicht-ärztliches
  Personal]{enumerate}{Assistenz}
  \item \StepHeading{Material-Vollständigkeit} prüfen
  \item \StepHeading{Sauerstoffberieselung} einschalten und auf
  15{\LPerMin*} einstellen, Reservoir des Beatmungsbeutels
  füllen
  \item Durchführendem melden, dass alle Materialien bereit sind
  \item \Z{Evtl. werden jetzt von einem Arzt Medikamente
  verabreicht, die die Intubation ermöglichen. Der
  Durchführende wird nun einige Zeit den Patienten mittels
  Beutel-/Masken-Beatmung beatmen (sofern nicht bereits
  geschehen).}
  \item \StepHeading{Laryngoskop zureichen}: in die \E{linke (!)}
  Hand zureichen\footnote{Es ist egal, ob der Durchführende
  Rechts- oder Linkshänder ist, das Laryngoskop wird
  \E{immer in der linken Hand} gehalten.}
  \item Auf Anweisung \emph{\enquote{Kehlkopfdruck},
  \enquote{Krikoiddruck}} auf den Kehlkopf drücken
  \item \StepHeading{Tubus zureichen}: in die \E{rechte (!)} Hand
  \item Auf Anweisung absaugen
  \item Beatmungsmaske vom Beatmungsbeutel trennen
  \item \Z{Durchführender intubiert und entfernt Führungsdraht (Mandrin)}
  \item \StepHeading{Cuffen}: Mit luftgefüllter Spritze Cuff
  aufblasen
  \item \StepHeading{Beatmungsbeutel} an Tubus anschließen
  \item \StepHeading{Stethoskop} dem Durchführenden in die Ohren
  klemmen
  \item \StepHeading{Lagekontrolle} mittels Auskulatation: Es werden
  3 Atemhübe verabreicht, wobei das Stethoskop über über dem Magen, linkem Lungenflügel und
  rechtem Lungenflügel positioniert wird. Es wird dabei
  überprüft ob der Tubus richtig liegt.
  \item Der Durchführende übernimmt die Beatmung.
  \item Einführen des \StepHeading{Beißkeils} oder des Guedel-Tubus
  \item \StepHeading{Fixierung} des Tubus und des Beißschutzes
  mittels Mullbinde, Klebestreifen oder speziellen
  Fixationsmaterialien
  \item \Z{Ggfs. Beatmungsgerät einstellen und
  anschließen}
\end{enumerate}
\end{ProcedureStepsWide}
% }{}{}{}{}


% \begin{PictureSeriesWide}{}{Bilderserie: Intubation \SourcePicture{Motal}}
% \PictureSeriesRowWide{4}
% {./images/motal-aass/intubation/dsc_4829-AASS-0030mm}
% {Zusammenbauen des Laryngoskops}
% {./images/motal-aass/intubation/dsc_4831-AASS-0030mm}
% {Einsprühen des Tubus \E{in der Verpackung}}
% {./images/motal-aass/intubation/dsc_4834-AASS-0030mm}
% {Präoxygenierung}
% {./images/motal-aass/intubation/dsc_4835-AASS-0030mm}
% {Zureichen des Laryngoskops}
% \PictureSeriesRowWide{4}
% {./images/motal-aass/intubation/dsc_4837-AASS-0030mm}
% {Zureichen des Tubus}
% {./images/motal-aass/intubation/dsc_4839-AASS-0030mm}
% {Einführen des Tubus, Assistenz führt auf Anweisung Kehlkopfdruck aus}
% {./images/motal-aass/intubation/dsc_4841-AASS-0030mm}
% {Cuffen nach Entfernung des Führungsdrahtes}
% {./images/motal-aass/intubation/dsc_4843-AASS-0030mm}
% {Stethoskop zureichen}
% \PictureSeriesRowWide{4}
% {./images/motal-aass/intubation/dsc_4844-AASS-0030mm}
% {Abhören des Patienten}
% {./images/motal-aass/intubation/dsc_4847-AASS-0030mm}
% {Beißkeil (hier: Guedel-Tubus) einführen}
% {./images/motal-aass/intubation/dsc_4849-AASS-0030mm}
% {Endgültige Fixierung des Tubus: Bis dahin muss der Tubus manuell fixiert werden!}
% {empty}
% {}
% \end{PictureSeriesWide}

\begin{PictureSeriesWide}{}{Bilderserie: \EEE{Intubation} \SourcePicture{Motal}}
\PictureSeriesRowWide{3}
{./images/motal-aass/intubation/dsc_4829-AASS-0030mm}
{Zusammenbauen des Laryngoskops}
{./images/motal-aass/intubation/dsc_4831-AASS-0030mm}
{Einsprühen des Tubus \E{in der Verpackung}}
{./images/motal-aass/intubation/dsc_4834-AASS-0030mm}
{Präoxygenierung}
{empty}
{}
\PictureSeriesRowWide{3}
{./images/motal-aass/intubation/dsc_4835-AASS-0030mm}
{Zureichen des Laryngoskops (\EE{linke} Hand!)}
{./images/motal-aass/intubation/dsc_4837-AASS-0030mm}
{Zureichen des Tubus (\EE{rechte} Hand!)}
{./images/motal-aass/intubation/dsc_4839-AASS-0030mm}
{Einführen des Tubus, Assistenz führt auf Anweisung Kehlkopfdruck aus}
{empty}
{}
\PictureSeriesRowWide{3}
{./images/motal-aass/intubation/dsc_4841-AASS-0030mm}
{Cuffen nach Entfernung des Führungsdrahtes}
{./images/motal-aass/intubation/dsc_4843-AASS-0030mm}
{Stethoskop zureichen}
{./images/motal-aass/intubation/dsc_4844-AASS-0030mm}
{Abhören des Patienten}
{empty}
{}
\PictureSeriesRowWide{3}
{./images/motal-aass/intubation/dsc_4847-AASS-0030mm}
{Beißkeil (hier: Guedel-Tubus) einführen}
{./images/motal-aass/intubation/dsc_4849-AASS-0030mm}
{Endgültige Fixierung des Tubus: Bis dahin muss der Tubus manuell fixiert werden!}
{empty}
{}
{empty}
{}
\end{PictureSeriesWide}

\normalfont

\begin{Literary}
  \fullcite{ArbeitstechnikenAZ1}

  \fullcite{InjektionenLeichtGemacht4}

  \fullcite{FamPropAeGruFert2011}
\end{Literary}



\nocite{ZechaOpGehilfe1}
\ChapterEnd